\begin{table}[tbp]
\centering\small
\caption{The Original Region with a Finite Surrender Charge}
\label{tab:r9_fee_certificate}
\begin{tabular}{lrrc}
\toprule
Contract & $10^4\Delta^{\rm adj}$ bound & $10^4\Delta^0$ bound & Both signs \\
\midrule
$m=8$ & [0.6869, 3.2473] & [-3.6968, -0.5280] & Yes \\
$F=0.80,\ m=1$ & [0.6869, 3.2473] & [-3.6968, 13.9258] & No \\
$F\in[0.85,0.90],\ m=1$ & [0.6869, 3.2473] & [-3.6968, -0.5280] & Yes \\
\bottomrule
\end{tabular}
\begin{minipage}{0.98\textwidth}\footnotesize\medskip
Each row uses the original rectangle $[0,0.25]\times[0.4,0.45]$, both adjustment regimes, and four law subcells. Chord and count are executed separately and give the same displayed bounds. Every unscaled difference includes $2\times10^{-7}$ arithmetic padding relative to the stored arrays. The fee interval uses a joint fee--benefit--law coefficient bound. The failed $F=0.80$ row and its adverse reoptimized corner are retained.
\end{minipage}
\end{table}
